% Part: incompleteness
% Chapter: arithmetization-syntax
% Section: proofs-in-ax

\documentclass[../../../include/open-logic-section]{subfiles}

\begin{document}

\olfileid{inc}{art}{pax}
\olsection{\usetoken{P}{derivation}ی اصل‌موضوعی}

\begin{explain}
برای حسابی‌سازی !!{derivation}sی اصل‌موضوعی باید !!{derivation}s را با اعداد
بازنمایی کنیم. چون !!{derivation}s صرفاً دنباله‌هایی از !!{formula}s
هستند، رویکرد آشکار آن است که هر !!{derivation} را با رمز دنبالهٔ رمزهای
!!{formula}sی آن رمزگذاری کنیم.
\end{explain}

\begin{defn}
اگر $\delta$ !!{derivation}ی اصل‌موضوعی متشکل از !!{formula}sی
$!A_1$،~\dots،~$!A_n$ باشد، آنگاه $\Gn{\delta}$ برابر است با
\[
\tuple{\Gn{!A_1}, \dots, \Gn{!A_n}}.
\]
\end{defn}

\begin{ex}
  !!{derivation} بسیار سادهٔ زیر را در نظر بگیرید:
  \begin{derivation}
  1. & $!B \lif (!B \lor !A)$ \\
  2. & $(!B \lif (!B \lor !A)) \lif (!A  \lif (!B \lif (!B \lor !A)))$\\
  3. & $!A  \lif (!B \lif (!B \lor !A))$
  \end{derivation}
  عدد گودل این !!{derivation} برابر خواهد بود با
  \begin{align*}
  \openTuple\,
    & \Gn{!B \lif (!B \lor !A)}, \\
    &\Gn{(!B \lif (!B \lor !A)) \lif (!A  \lif (!B \lif (!B \lor !A)))},\\
    & \Gn{!A  \lif (!B \lif (!B \lor !A))} \,\closeTuple.
  \end{align*}
\end{ex}

\begin{explain}
پس از تثبیت بازنمایی !!{derivation}s باید نشان دهیم می‌توانیم چنین
!!{derivation}sیی را به‌طور بازگشتی اولیه دست‌کاری کنیم و ویژگی‌ها و
روابط اساسی‌شان را نیز چنین بیان کنیم. بعضی عملیات ساده‌اند: برای
نمونه، با داشتن عدد گودل~$d$ از !!a{derivation}ی،
$(d)_{\len{d}-1}$ عدد گودل !!{formula} پایانی آن را به دست می‌دهد.
بعضی موارد بسیار دشوارترند و دست‌کم طرحی از چگونگی انجامشان می‌دهیم.
هدف نشان‌دادن این است که رابطهٔ «$\delta$ !!a{derivation}ی از~$!A$ بر
پایهٔ~$\Gamma$ است» نسبت به اعداد گودل~$\delta$ و~$!A$ بازگشتی
اولیه است.
\end{explain}

\begin{prop}
\ollabel{prop:followsby}
روابط زیر بازگشتی اولیه‌اند:
\begin{enumerate}
\item $!A$ یک اصل موضوعه است.
\item سطر $i$اُم در~$\delta$ با وضع مقدم توجیه می‌شود.
\item سطر $i$اُم در~$\delta$ با~\QR{} توجیه می‌شود.
\item $\delta$ یک !!{derivation} درست است.
\end{enumerate}
\end{prop}

\begin{proof}
باید نشان دهیم روابط متناظر میان اعداد گودل !!{formula}s و اعداد
گودل !!{derivation}s بازگشتی اولیه‌اند.
\begin{enumerate}
\item فهرستی معین از شِماهای اصل موضوعه داریم و~$!A$ هنگامی اصل
  موضوعه است که شکلی مطابق یکی از این شِماها داشته باشد. چون فهرست
  شِماها متناهی است، کافی است نشان دهیم برای هر شِمای اصل موضوعه
  می‌توانیم به‌طور بازگشتی اولیه بیازماییم که آیا~$!A$ آن شکل را
  دارد. برای نمونه، شِمای اصل موضوعهٔ زیر را در نظر بگیرید:
  \[
  !B \lif (!C \lif !B).
  \]
  $!A$ نمونه‌ای از این شِماست اگر !!{formula}sیی چون~$!B$ و~$!C$
  وجود داشته باشند که از الحاق «$($»، سپس~$!B$، سپس «$\lif$»، سپس
  «$($»، سپس~$!C$، سپس «$\lif$»، سپس~$!B$ و سرانجام «$))$»،
  فرمول~$!A$ به دست آید. می‌توانیم ویژگی متناظرِ عدد گودل~$n$ برای
  $!A$ را بیازماییم، زیرا الحاق دنباله‌ها بازگشتی اولیه است و اعداد
  گودل~$!B$ و~$!C$ باید از عدد گودل~$!A$ کوچک‌تر باشند؛ زیرا هرگاه
  رابطه برقرار باشد، هر دو زیر!!{formula}s~$!A$ هستند. بنابراین می‌توانیم
  تعریف کنیم:
  \begin{multline*}
  \fn{IsAx}_{!B \lif (!C \lif !B)}(n) \defiff \bexists{b< n}{
    \bexists{c < n}{(\fn{Sent}(b) \land \fn{Sent}(c) \land {}}}\\ n =
  \Gn{(} \concat b \concat \Gn{\lif} \concat \Gn{(} \concat c \concat
  \Gn{\lif} \concat b \concat \Gn{))}).
  \end{multline*}
  اگر برای هر شِمای اصل موضوعه چنین تعریفی داشته باشیم، فصلِ
  آن‌ها ویژگی~$\fn{IsAx}(n)$، یعنی «$n$ عدد گودل یک اصل موضوعه
  است»، را تعریف می‌کند.
\item سطر $i$اُم در~$\delta$ با وضع مقدم توجیه می‌شود اگر و تنها اگر
  سطرهایی چون~$j$ و~$k<i$ وجود داشته باشند که !!{sentence}ٔ سطر~$j$
  فرمولی چون~$!A$، جملهٔ سطر~$k$ برابر~$!A\lif !B$، و جملهٔ
  سطر~$i$ برابر~$!B$ باشد:
  \begin{multline*}
    \fn{MP}(d, i) \defiff \bexists{j < i}{\bexists{k < i}{}}\\
    (d)_k = \Gn{(} \concat (d)_j
      \concat \Gn{\lif} \concat (d)_i \concat \Gn{)}.
  \end{multline*}
  چون کمیت‌گذاری کراندار، الحاق و~$=$ بازگشتی اولیه‌اند، این تعریف
  رابطه‌ای بازگشتی اولیه به دست می‌دهد.
\item سطری در~$\delta$ با~\QR{} توجیه می‌شود اگر به صورت
  $!B\lif\lforall[x][!A(x)]$ باشد، سطری پیشین برای بعضی !!{constant}~$c$
  برابر~$!B\lif !A(c)$ باشد، و~$c$ در~$!B$ رخ ندهد. این درست است اگر
  و تنها اگر
  \begin{enumerate}
  \item !!a{sentence}ای چون~$!B$ و
  \item !!a{formula}ی چون~$!A(x)$ با تنها متغیر آزاد~$x$ وجود داشته
    باشند به‌گونه‌ای که
  \item سطر~$i$ شامل~$!B\lif\lforall[x][!A(x)]$ باشد،
  \item سطری چون~$j<i$ شامل~$!B\lif\Subst{!A}{c}{x}$ برای
    ثابتی چون~$c$ باشد،
  \item و این~$c$ در~$!B$ رخ ندهد.
  \end{enumerate}
  همهٔ این‌ها را می‌توان به‌طور بازگشتی اولیه آزمود، زیرا اعداد گودل
  $!B$، $!A(x)$ و~$x$ از عدد گودل فرمول سطر~$i$، و عدد گودل~$c$ از
  عدد گودل فرمول سطر~$j$ کوچک‌ترند:
  \begin{multline*}
    \fn{QR}_1(d, i) \defiff \bexists{j<i}{\bexists{b < (d)_i}{\bexists{x <
        (d)_i}{\bexists{a < (d)_i}{\bexists{c < (d)_j}{(}}}}}
    \\ \fn{Var}(x) \land \fn{Const}(c) \land {} \\
    (d)_i = \Gn{(} \concat
    b \concat \Gn{\lif} \concat \Gn{\lforall} \concat x \concat a
    \concat \Gn{)} \land {}\\
    (d)_j = \Gn{(} \concat b \concat
    \Gn{\lif} \concat \fn{Subst}(a,c,x) \concat \Gn{)} \land {}\\ \fn{Sent}(b)
    \land \fn{Sent}(\fn{Subst}(a,c,x)) \land {} \bforall{k <
      \len{b}}{(b)_k \neq (c)_0}).
  \end{multline*}
  اینجا فرض می‌کنیم $c$ و~$x$ اعداد گودل ثابت و متغیرِ درنظرگرفته‌شده
  به‌عنوان ترم‌اند، نه رمز نمادشان. اینکه $x$ تنها متغیر آزادِ
  $!A(x)$ است با آزمودن اینکه~$\Subst{!A(x)}{c}{x}$ یک !!a{sentence} است بررسی
  می‌شود؛ و با الزام متفاوت‌بودن هر نماد~$!B$ از~$c$ تضمین می‌کنیم
  که~$c$ در~$!B$ رخ ندهد.

  صورت دیگر~\QR{} را به‌عنوان تمرین وا می‌گذاریم.
\item $d$ عدد گودل یک !!{derivation} درست است اگر و تنها اگر هر سطر آن
  اصلی موضوعه باشد یا با وضع مقدم یا~\QR{} توجیه شود. بنابراین:
  \[
  \fn{Deriv}(d) \defiff \bforall{i < \len{d}}{(\fn{IsAx}((d)_i) \lor
    \fn{MP}(d,i) \lor \fn{QR}(d, i))}.
  \]
\end{enumerate}
\end{proof}

\begin{prob}
روابط زیر را مانند \olref[inc][art][pax]{prop:followsby} تعریف کنید:
\begin{enumerate}
\item $\fn{IsAx}_{!A \lif (!B \lif (!A \land !B))}(n)$؛
\item $\fn{IsAx}_{\lforall[x][!A(x)] \lif !A(t)}(n)$؛
\item $\fn{QR}_{2}(d, i)$ (برای صورت دیگر~\QR).
\end{enumerate}
\end{prob}

\begin{prop}
فرض کنید $\Gamma$ مجموعه‌ای بازگشتی اولیه از !!{sentence}s باشد. آنگاه
رابطهٔ~$\Prf[\Gamma](x,y)$ که بیان می‌کند «$x$ رمز
!!a{derivation}~$\delta$ از~$!A$ بر پایهٔ~$\Gamma$ است و~$y$ عدد گودل
$!A$ است»، بازگشتی اولیه است.
\end{prop}

\begin{proof}
فرض کنید «$y\in\Gamma$» با محمول بازگشتی اولیهٔ~$R_\Gamma(y)$ داده
شود. باید نشان دهیم رابطهٔ~$\Prf[\Gamma](x,y)$ بازگشتی اولیه است؛
$\Prf[\Gamma](x,y)$ برقرار است اگر و تنها اگر $y$ عدد گودل
!!a{sentence}ای چون~$!A$
و~$x$ رمز !!a{derivation}ی از~$!A$ بر پایهٔ~$\Gamma$ باشد.

بنا بر گزارهٔ پیشین، ویژگی~$\fn{Deriv}(x)$ که برقرار است اگر و تنها
اگر $x$ رمز !!{derivation} درست~$\delta$ باشد، بازگشتی اولیه است. اما آن
تعریف، مجموعهٔ~$\Gamma$ را به‌عنوان راهی دیگر برای توجیه سطرهای
!!{derivation} در نظر نمی‌گرفت. آزمون بازگشتی اولیهٔ ما برای توجیه‌شدن
سطری با~\QR{} نیز این شرط را نادیده گرفت که ثابت~$c$ مجاز نیست
در~$\Gamma$ رخ دهد. می‌توان تعریف را چنان اصلاح کرد که~$\Gamma$ را
مستقیماً در نظر بگیرد، اما استفاده از~$\fn{Deriv}$ و قضیهٔ استنتاج
آسان‌تر است. $\Gamma\Proves !A$ اگر و تنها اگر فهرستی متناهی از
!!{sentence}sی $!B_1$،~\dots،~$!B_n\in\Gamma$ وجود داشته باشد که
$\{!B_1,\dots,!B_n\}\Proves !A$. و بنا بر قضیهٔ استنتاج، این درست
است هرگاه
$\Proves(!B_1\lif(!B_2\lif\cdots(!B_n\lif !A)\cdots))$.
می‌توان به‌طور بازگشتی اولیه آزمود که آیا !!a{sentence}ای با عدد گودل~$z$
این شکل را دارد. پس به‌جای آنکه $x$ را عدد گودل !!a{derivation} !!{sentence}ای
با عدد گودل~$y$ \emph{بر پایهٔ~$\Gamma$} بدانیم، آن را عدد گودل
!!a{derivation} یک شرطی تودرتو با شکل بالا بر پایهٔ~$\emptyset$ می‌گیریم.

نخست، اگر دنباله‌ای از !!{sentence}s داشته باشیم، می‌توانیم به‌طور
بازگشتی اولیه شرطی‌ای بسازیم که همهٔ این جمله‌ها مقدم و !!{sentence}ای
معین تالی آن باشند:
\begin{align*}
  \fn{hCond}(s, y, 0) & = y \\
  \fn{hCond}(s, y, n+1) & = \Gn{(} \concat (s)_{n} \concat \Gn{\lif}
  \concat \fn{hCond}(s, y, n) \concat \Gn{)}\\
  \fn{Cond}(s, y) & = \fn{hCond}(s, y, \len{s})\\
  \intertext{پس می‌توانیم $\Prf[\Gamma](x,y)$ را چنین تعریف کنیم:}
  \Prf[\Gamma](x, y) & \defiff \bexists{s < \fn{sequenceBound}(x,x)}{(} \\
  &\qquad (x)_{\len{x}-1} = \fn{Cond}(s,y) \land {} \\
  &\qquad\bforall{i<\len{s}}{(s)_i \in \Gamma} \land {} \\
  &\qquad\fn{Deriv}(x)).
\end{align*}
کران~$s$ از این ملاحظه به دست می‌آید که هر~$(s)_i$ عدد گودل یک
زیر!!{formula} از آخرین سطر !!{derivation} است و بنابراین از
$(x)_{\len{x}-1}$ کوچک‌تر است. تعداد مقدم‌های~$!B\in\Gamma$، یعنی
طول~$s$، نیز از طول آخرین سطر~$x$ کوچک‌تر است.
\end{proof}

\end{document}
